\documentclass[11pt,a4paper]{article}
\usepackage[margin=1in]{geometry}
\usepackage[T1]{fontenc}
\usepackage[utf8]{inputenc}
\usepackage{amsmath,amssymb}
\usepackage{hyperref}
\usepackage{setspace}
\usepackage{enumitem}
\hypersetup{colorlinks=true,linkcolor=black,urlcolor=blue}
\onehalfspacing

\begin{document}

\begin{center}
{\large\bfseries Cover Letter}\\[6pt]
{\normalsize Submission to \emph{Computers and Electronics in Agriculture}}
\end{center}

\vspace{1.2em}

\noindent\textbf{Date:} October 14, 2026

\noindent\textbf{To:} The Editor, \emph{Computers and Electronics in Agriculture} (Elsevier)

\noindent\textbf{Manuscript title:} ``Fine-Tuning or Frozen Probes? Empirical Applicability Boundaries for Off-the-Shelf CLIP in Tobacco Leaf-Disorder Diagnosis''

\noindent\textbf{Authors:} Chenfeng Long$^{1,2,3}$, Zhiyong Tan$^{1,2,3}$, Yangjun Deng$^{1,2,3}$, Xiang Luo$^{4}$, Lizhuo Zhang$^{1,2,3,*}$

\begin{itemize}[leftmargin=2em,itemsep=2pt,topsep=4pt]
\item[$^{1}$] College of Information and Intelligence, Hunan Agricultural University, Changsha 410128, China
\item[$^{2}$] Yuelushan Laboratory, Changsha 410128, China
\item[$^{3}$] Hunan Provincial Engineering and Technology Research Center for Rural and Agricultural Informatization, Hunan Agricultural University, Changsha 410128, China
\item[$^{4}$] Information Department, China Guangdong Tobacco Meizhou Co.\ Ltd, Meizhou 514000, China
\item[$^{*}$] Corresponding author: zhanglizhuo@hunau.edu.cn
\end{itemize}

\vspace{1.5em}

Dear Editor,

We are pleased to submit the above manuscript for consideration as a Research Article in \emph{Computers and Electronics in Agriculture}.

\medskip

\noindent\textbf{Why this work matters for your readership.}
Whether off-the-shelf vision--language models can classify fine-grained crop diseases without domain-specific training is a practical question for the resource-constrained deployments that \emph{Computers and Electronics in Agriculture} addresses. On a primary 16-class tobacco leaf-disorder dataset (TPDD), with cross-crop checks on tomato and corn (PlantVillage) and a field-condition evaluation, we report findings of direct practical consequence. First, generic CLIP is unreliable zero-shot for this fine-grained task: on the two tobacco sets no prompt template is detectably better than the majority-class baseline, and knowledge-graph-augmented prompts give no reliable gain. Second, the CLIP-versus-ResNet ordering varies with the adaptation protocol rather than the data regime: with \emph{frozen} CLIP probes at the default $C{=}1$, the apparent few-shot CLIP advantage over the legacy (under-tuned) CNN is an under-tuning artifact, and against a matched-budget fine-tuned ResNet-18 the frozen probe ties at $K{=}1$ and trails at $K{\geq}2$; whereas with \emph{both backbones fine-tuned} under the same matched budget, fine-tuned CLIP leads fine-tuned ResNet-18 at every $K$. Third, at full data fine-tuned ResNet-18 holds a small, sign-stable edge over the validation-regularized frozen CLIP probe under 5-fold cross-validation. Together these give practitioners an empirical boundary for when off-the-shelf CLIP is and is not a reasonable, training-light option for fine-grained leaf-disorder diagnosis.

\medskip

\noindent\textbf{Fit and originality.}
The work is empirical and methods-comparative, in line with the journal's scope on intelligent and electronic systems for agriculture. Its contribution is a within-crop empirical boundary, tested under seed-paired statistics and multiplicity correction, with a cross-crop replication---rather than a new architecture. We deliberately frame the contribution as protocol-controlled boundary conditions for deployment, not as a new model: the CLIP-based plant-disease numbers in the current literature are not protocol-comparable (tuning budgets, evaluation splits and class compositions differ across reports), and practitioners lack an evidence base for deciding between frozen-probe and fine-tuning workflows on fine-grained leaf disorders. We believe this is precisely the kind of applied decision evidence \emph{Computers and Electronics in Agriculture} readers need.

\medskip

\noindent\textbf{Declarations.}
\begin{enumerate}[leftmargin=2.5em,itemsep=3pt,topsep=4pt]
\item This manuscript is original, has not been published, and is not under consideration by any other journal or conference.
\item All authors have read and approved the submission. X.~Luo is employed by China Guangdong Tobacco Meizhou Ltd, and part of this work was supported by the Meizhou Tobacco Science Research Project (Grant~202404); the authors declare no other conflicts of interest.
\item AI-assisted tools (an LLM coding assistant) were used for drafting and code assistance during preparation; the authors critically reviewed every claim, verified all reported statistics against the underlying result files, and take full responsibility for the content.
\item The code, per-seed result files, per-split label manifests, and audit files supporting all findings are available in the public repository \url{https://github.com/zhiyongtanzhiyong-svg/probe-protocol} so that the reported statistics can be independently re-verified during review; the precomputed feature stores ($\approx$2\,GB) and the re-derived TPDD image set (the original 11\,MB release is no longer available to us) are available from the authors on request.
\end{enumerate}

\medskip

Thank you for your consideration. We look forward to your response.

\vspace{1.5em}

Sincerely,

\vspace{1.2em}

Lizhuo Zhang (on behalf of all co-authors)\\
College of Information and Intelligence, Hunan Agricultural University\\
Changsha 410128, China\\
\texttt{zhanglizhuo@hunau.edu.cn}

\end{document}
